\section*{Chapter \ref{ch:g11:dna-replication} --- DNA Replication}

\begin{solution}{exo:g11:dna-replication:1}
Each strand serves as a template on which a new complementary strand
is built. It works only because A pairs with T and G with C: the
sequence of one strand fixes that of the other, so a strand alone
carries the whole information.
\end{solution}

\begin{solution}{exo:g11:dna-replication:2}
Semi-conservative: each daughter molecule keeps one parent strand and
gains one new strand. Conservative: the parent molecule stays whole
and a completely new double strand is built beside it.
\end{solution}

\begin{solution}{exo:g11:dna-replication:3}
The S phase of the \omterm{def:g10:cells-common-unit:cell}{cell} cycle; the amount of \omterm{def:g10:universal-dna:information}{DNA} per \omterm{def:g10:cells-common-unit:organelle}{nucleus} doubles.
\end{solution}

\begin{solution}{exo:g11:dna-replication:4}
\omterm{prop:g11:dna-replication:forks}{DNA polymerase} adds to a growing strand the \omterm{def:g10:universal-dna:nucleotide}{nucleotides} complementary
to the template, one at a time. A \omterm{prop:g11:dna-replication:forks}{replication fork} is the moving point
where the parent strands are separated and the two new strands are
being built.
\end{solution}

\begin{solution}{exo:g11:dna-replication:5}
\texttt{CTAATGTCCG}, base under base.
\end{solution}

\begin{solution}{exo:g11:dna-replication:6}
Conservative, generation 1: a heavy band and a light band, no hybrid
--- but only a hybrid band was seen. Dispersive, generation 2: a
single band a quarter of the way from light to heavy --- but two
distinct bands were seen. Each model predicts bands that were not
observed.
\end{solution}

\begin{solution}{exo:g11:dna-replication:7}
$2^5 = 32$ molecules, 2 hybrid and 30 light: a light band fifteen
times the hybrid one, no heavy band.
\end{solution}

\begin{solution}{exo:g11:dna-replication:8}
Generation 0: one light band. Generation 1: one hybrid band.
Generation 2: hybrid and heavy bands of equal intensity --- the mirror
image of the original experiment.
\end{solution}

\begin{solution}{exo:g11:dna-replication:9}
$\num{2.5e8}/50 = \num{5e6}\,\unit{s}$, about 58 days --- far longer
than 8 hours. The \omterm{def:g10:universal-dna:information}{chromosome} must be copied from many origins at once,
each opening two forks.
\end{solution}

\begin{solution}{exo:g11:dna-replication:10}
Both strands of $\num{3.2e9}$ base pairs, twice: $\num{6.4e9}$ base
pairs, i.e. $\num{6.4e9}$ new \omterm{def:g10:universal-dna:nucleotide}{nucleotides} joined (one per base pair
copied). At $10^{-9}$ per \omterm{def:g10:universal-dna:nucleotide}{nucleotide}, about 6 errors.
\end{solution}

\begin{solution}{exo:g11:dna-replication:11}
Opening a bubble exposes template on both sides of the origin; each
fork copies away from the origin, so the two go in opposite directions
and together cover the whole region. When two bubbles meet, their
forks fuse and the daughter strands join end to end into continuous
molecules.
\end{solution}

\begin{solution}{exo:g11:dna-replication:12}
The bacteria stop dividing (no replication, no division). The wound
stops healing, since its repair needs \omterm{def:g10:cells-common-unit:cell}{cell} divisions. The nerve \omterm{def:g10:cells-common-unit:cell}{cell}
is unaffected: it no longer replicates its \omterm{def:g10:universal-dna:information}{DNA}. A drug that stops
rapidly dividing \omterm{def:g10:cells-common-unit:cell}{cells} while sparing non-dividing ones suggests a
treatment against bacteria or against dividing cancer \omterm{def:g10:cells-common-unit:cell}{cells}.
\end{solution}

\begin{solution}{exo:g11:dna-replication:13}
The two original strands are never destroyed: at every generation each
is again paired with a new light strand, so exactly two hybrid
molecules always exist. After 10 generations: $2/2^{10} = 2/1024
\approx 0.2\%$ --- present but too faint to see.
\end{solution}

\begin{solution}{exo:g11:dna-replication:14}
$\num{6.4e9} \times 10^{-5} = \num{64000}$ errors per division. With
twenty thousand \omterm{def:g10:universal-dna:gene}{genes}, hundreds of \omterm{def:g10:universal-dna:gene}{genes} would be damaged at every
division; after a few divisions the \omterm{def:g10:cells-common-unit:cell}{cells} would carry lethal defects
and the line would die out.
\end{solution}

\begin{solution}{exo:g11:dna-replication:15}
A new round of replication starts at the origin before the previous
round has finished: the \omterm{def:g10:universal-dna:information}{chromosome} carries several nested pairs of
forks at once. Each division then receives a \omterm{def:g10:universal-dna:information}{chromosome} already partly
copied for the next, and divisions can follow one another faster than
a single copying takes.
\end{solution}

\begin{solution}{pb:g11:dna-replication:1}
\textbf{1.} Nitrogen is 16\% of the mass and is 7\% heavier:
$0.16 \times 0.07 \approx 1.1\%$. The gradient method separates
densities differing by a fraction of a per cent, so 1\% gives
well-separated bands.

\textbf{2.} So that essentially all the \omterm{def:g10:universal-dna:information}{DNA}, not just half or three
quarters, was heavy: after 14 generations less than one part in
$10^4$ of the original light strands remains.

\textbf{3.} Exactly half-way between the two bands: its density is the
average.

\textbf{4.} A single bacterium holds a few femtograms of \omterm{def:g10:universal-dna:information}{DNA}; the
photograph needs micrograms, hence billions of \omterm{def:g10:cells-common-unit:cell}{cells}.

\textbf{5.} By keeping a culture in heavy medium as a control and
checking that its band stayed heavy, and by checking the light band
of bacteria grown in light medium throughout.

\textbf{6.} Semi-conservative: generation 1, one hybrid band;
generation 2, hybrid and light bands, equal.

\textbf{7.} Conservative: generation 1, heavy and light bands, equal;
generation 2, one heavy for three light.

\textbf{8.} Dispersive: generation 1, one band at the hybrid position;
generation 2, one band a quarter of the way from light to heavy.

\textbf{9.} The conservative model, which predicted no hybrid band, is
eliminated. Semi-conservative and dispersive both predict the single
hybrid band and remain.

\textbf{10.} Semi-conservative survives: it alone predicts two
separate bands. The dispersive model predicts a single band drifting
lighter, never splitting.

\textbf{11.} $2^n$ molecules, of which exactly 2 carry an original
strand.

\textbf{12.} Light : hybrid $= (2^n - 2) : 2$: generation 3, $3:1$;
generation 4, $7:1$; generation 6, $31:1$.

\textbf{13.} Hybrid fraction $2/2^n < 0.05$ when $2^n > 40$: from
generation 6 ($2/64 \approx 3\%$).

\textbf{14.} Two bands, one heavy and one light, of equal amount: a
hybrid molecule is one heavy strand plus one light strand. This tests
the dispersive model, which predicts every single strand to be of
intermediate density, and eliminates it independently.

\textbf{15.} Starting light, one generation in heavy medium gives one
hybrid band and two generations hybrid plus heavy: the same pattern
mirrored, showing the result does not depend on which isotope is
"old".

\textbf{16.} $\num{4.6e6}/(2 \times 1000) = \qty{2300}{s}$, about
38 minutes.

\textbf{17.} 77 minutes. On a circle the two strands are exposed on
both sides of any opening; two forks leaving one origin in opposite
directions is the simplest way to cover the whole circle.

\textbf{18.} One origin copies $2 \times 50 \times 8 \times 3600 =
\num{2.88e6}$ base pairs in 8 hours; $\num{6.4e9}/\num{2.88e6}
\approx 2200$ origins at the very least (in reality tens of thousands,
firing at different times).

\textbf{19.} Bacterium: $\num{9.2e6} \times 10^{-9} \approx 0.01$ error
per replication; human \omterm{def:g10:cells-common-unit:cell}{cell}: about 6. Most of the human genome is not
\omterm{def:g10:universal-dna:gene}{genes}, so most errors fall where they change nothing, and the \omterm{def:g10:cells-common-unit:cell}{cell} has
two copies of every \omterm{def:g10:universal-dna:gene}{gene}.

\textbf{20.} The two bands at generation two --- hybrid and light in
equal amounts --- proved replication semi-conservative; the two forks
copy the bacterial \omterm{def:g10:universal-dna:information}{chromosome} in about 38 minutes, roughly the
bacterium's generation time, so copying paces division.
\end{solution}
